% CP-VI-0017
% target_chapter: VI/15 — Kernels, Images, and Quotients
% source_unit_id: IMP-DL-0E3E7EC766-P017
% source: Downloads/theory-of-algebraic-geometry-1.tex L3489-L3497
% solution_provenance: SOURCE_EXPLICIT_SOLUTION

\subsection*{Exercise 5}

Show that if
\[
f : \mathcal F \longrightarrow \mathcal G
\]
is a morphism between sheaves, then the sheaf image $\operatorname{im} f$
can be naturally identified with a subsheaf of $\mathcal G$.

\subsection*{Solution}

Let
\[
f : \mathcal F \longrightarrow \mathcal G
\]
be a morphism of sheaves of abelian groups on $X$.

For every open set $U \subseteq X$, we have a group homomorphism
\[
f_U : \mathcal F(U) \longrightarrow \mathcal G(U).
\]

One might try to define the image presheaf by
\[
U \longmapsto \operatorname{im}(f_U).
\]
This is indeed a presheaf, but it need not be a sheaf.

The reason is that being in the image of $f_U$ is a global condition.
A section of $\mathcal G(U)$ might be locally in the image of $f$, but fail to
come from a single global section of $\mathcal F(U)$.

Therefore the correct sheaf-theoretic image is the sheafification of the
presheaf image.

Define the presheaf image
\[
\operatorname{Im}^{\operatorname{pre}}(f)
\]
by
\[
\operatorname{Im}^{\operatorname{pre}}(f)(U)
=
\operatorname{im}
\left(
f_U : \mathcal F(U) \to \mathcal G(U)
\right).
\]

Then the sheaf image is defined as
\[
\operatorname{im} f
:=
\left(\operatorname{Im}^{\operatorname{pre}}(f)\right)^+.
\]

We claim that $\operatorname{im} f$ can be naturally identified with a subsheaf
of $\mathcal G$.

\subsection*{1. Explicit description of the image sheaf}

Define a subobject
\[
\mathcal H \subseteq \mathcal G
\]
as follows.

For every open set $U \subseteq X$, let
\[
\mathcal H(U)
=
\left\{
s \in \mathcal G(U)
\ \middle|\
\begin{array}{l}
	\text{for every } p \in U, \text{ there exists an open neighborhood}
	\\
	V \subseteq U \text{ of } p
	\text{ and } t \in \mathcal F(V)
	\text{ such that}
	\\
	s|_V = f_V(t)
\end{array}
\right\}.
\]

In words:

\[
\mathcal H(U)
=
\{\text{sections of } \mathcal G(U)
\text{ which are locally in the image of } f\}.
\]

Clearly
\[
\mathcal H(U) \subseteq \mathcal G(U).
\]

We now show that $\mathcal H$ is a subsheaf of $\mathcal G$.

\subsection*{2. $\mathcal H$ is a subpresheaf of $\mathcal G$}

Let
\[
s \in \mathcal H(U)
\]
and let
\[
W \subseteq U
\]
be open.

We must show that
\[
s|_W \in \mathcal H(W).
\]

Take any point $p \in W$.
Since $s \in \mathcal H(U)$, there is an open neighborhood
\[
V \subseteq U
\]
of $p$ and a section
\[
t \in \mathcal F(V)
\]
such that
\[
s|_V = f_V(t).
\]

Then
\[
V \cap W
\]
is an open neighborhood of $p$ in $W$, and
\[
(s|_W)|_{V \cap W}
=
s|_{V \cap W}
=
f_{V \cap W}(t|_{V \cap W}).
\]

Thus $s|_W$ is locally in the image of $f$.

Therefore
\[
s|_W \in \mathcal H(W).
\]

So $\mathcal H$ is a subpresheaf of $\mathcal G$.

\subsection*{3. $\mathcal H$ is a sheaf}

Let $U \subseteq X$ be open and let $\{U_i\}_{i \in I}$ be an open cover of $U$.

Suppose that
\[
s_i \in \mathcal H(U_i)
\]
and
\[
s_i|_{U_i \cap U_j}
=
s_j|_{U_i \cap U_j}
\]
for all $i,j$.

Since $\mathcal G$ is a sheaf, there exists a unique section
\[
s \in \mathcal G(U)
\]
such that
\[
s|_{U_i}=s_i
\]
for every $i$.

We must show that
\[
s \in \mathcal H(U).
\]

Let $p \in U$.
Choose $i$ such that
\[
p \in U_i.
\]

Since
\[
s_i \in \mathcal H(U_i),
\]
there is an open neighborhood
\[
V \subseteq U_i
\]
of $p$ and a section
\[
t \in \mathcal F(V)
\]
such that
\[
s_i|_V = f_V(t).
\]

But
\[
s|_V = s_i|_V.
\]
Therefore
\[
s|_V = f_V(t).
\]

Thus $s$ is locally in the image of $f$ at every point $p \in U$.

Hence
\[
s \in \mathcal H(U).
\]

So $\mathcal H$ is a sheaf.

Since each $\mathcal H(U)$ is a subgroup of $\mathcal G(U)$ and restrictions
are inherited from $\mathcal G$, $\mathcal H$ is a subsheaf of $\mathcal G$.

\subsection*{4. Identification with the sheaf image}

We now show that
\[
\mathcal H \cong \operatorname{im} f.
\]

The presheaf image
\[
\operatorname{Im}^{\operatorname{pre}}(f)
\]
is a subpresheaf of $\mathcal H$, because if
\[
s = f_U(t)
\]
globally on $U$, then certainly $s$ is locally in the image of $f$.

Thus we have a morphism of presheaves
\[
\operatorname{Im}^{\operatorname{pre}}(f)
\longrightarrow
\mathcal H.
\]

Since $\mathcal H$ is a sheaf, the universal property of sheafification gives
a unique morphism of sheaves
\[
\left(\operatorname{Im}^{\operatorname{pre}}(f)\right)^+
\longrightarrow
\mathcal H.
\]

That is,
\[
\operatorname{im} f \longrightarrow \mathcal H.
\]

Conversely, every section of $\mathcal H$ is locally in the presheaf image.
Therefore it is exactly a section obtained by gluing local sections of the
presheaf image. This is precisely what sheafification does.

Hence
\[
\operatorname{im} f \cong \mathcal H.
\]

Therefore
\[
\boxed{
	\operatorname{im} f(U)
	=
	\left\{
	s \in \mathcal G(U)
	\mid
	s \text{ is locally of the form } f(t)
	\right\}
}
\]
and this realizes $\operatorname{im} f$ as a subsheaf of $\mathcal G$.

\subsection*{5. Stalkwise interpretation}

For every point $p \in X$,
\[
(\operatorname{im} f)_p
=
\operatorname{im}(f_p),
\]
where
\[
f_p : \mathcal F_p \longrightarrow \mathcal G_p.
\]

Indeed, taking stalks commutes with sheafification and with images in this
situation.

Thus a germ
\[
s_p \in \mathcal G_p
\]
lies in $(\operatorname{im} f)_p$ precisely when there exists a germ
\[
t_p \in \mathcal F_p
\]
such that
\[
s_p = f_p(t_p).
\]

So the sheaf image is the subsheaf of $\mathcal G$ whose sections are locally
images of sections of $\mathcal F$.

\[
\boxed{
	\operatorname{im} f
	\subseteq
	\mathcal G
}
\]

\newpage
